% The reader's guide sat inside the old single-file front matter. Under the
% template it becomes its own included file, placed after the acronym table
% because it expands on the same vocabulary at greater length.

\clearpage
{
    \pagestyle{guide}
    \chapter*{Reader's Guide: Modes, Codes and Terms}
    \addcontentsline{toc}{chapter}{Reader's Guide: Modes, Codes and Terms}

This dissertation compares six ways of showing a web page to the same program,
across eight combinations of website and model. Those six ways and eight
combinations are referred to by short codes throughout. This page decodes them,
and defines the recurring terms in plain language. Nothing here is needed to
follow the argument of Chapter~\ref{ch:intro}; it is placed first so that the
tables and figures in later chapters can be read without hunting for a
definition.

\paragraph{The six observation modes.} Each mode is one complete recipe for what
the program is shown at a single browser step. They differ in three things: the
page text supplied, the wording of the instructions around that text, and
whether an image is sent.

\begin{center}\small
\begin{tabularx}{\textwidth}{@{}lX@{}}
\toprule
Code & What the program is shown \\
\midrule
\mDOM{}  & A structured text listing of the page's elements, their labels and
           their roles. No image. \\
\mPP{}   & The same text listing, but with instructions written as though a
           numbered screenshot were coming. No image. \\
\mPT{}   & Only the lines of that listing that carry an element number, with
           the structure flattened away. No image. \\
\mPS{}   & Those same numbered lines, with instructions written as though a
           numbered screenshot were coming. No image. \\
\mSoM{}  & Those same numbered lines \emph{and} a screenshot with the same
elements numbered and boxed on it. \\
\mVis{}  & The plain screenshot only. No page text at all. \\
\bottomrule
\end{tabularx}
\end{center}

\noindent The four modes with no image are the cheap end; \mSoM{} is the
expensive end. \mPT{}, \mPP{} and \mPS{} exist to separate three changes that
happen together when one moves from \mDOM{} to \mSoM{}, and are called
\emph{phantom} modes throughout (Section~\ref{sec:modes}).

\paragraph{The eight website-and-model combinations.} A \emph{cell} is one
website paired with one model. Every cell runs all six modes over the same
tasks.

\begin{center}\small
\begin{tabularx}{\textwidth}{@{}llX@{}}
\toprule
Code & Stands for & Which is \\
\midrule
\cell{cls}    & VisualWebArena classifieds & a self-hosted classified-ads site \\
\cell{red}    & VisualWebArena reddit      & a self-hosted discussion-forum site \\
\cell{wa\_red} & WebArena reddit           & the same forum in a second, text-oriented benchmark \\
\midrule
\cell{B0} & \texttt{Qwen3-VL-235B-A22B} & the large model, reached through a paid API \\
\cell{B1} & \texttt{Qwen3-VL-4B}        & a small model of the same design family, run locally \\
\cell{B2} & \texttt{google/gemma-3-4b-it} & a small model of a \emph{different} design family, run locally \\
\bottomrule
\end{tabularx}
\end{center}

\noindent So \cell{cls$\cdot$B0} means ``the classified-ads site, run with the
large model''. The eight cells are the three sites crossed with the models
available on each.

\paragraph{Recurring terms.}

\begin{center}\small
\begin{tabularx}{\textwidth}{@{}lX@{}}
\toprule
Term & Meaning in this dissertation \\
\midrule
task success & Whether the benchmark's own automatic checker judged the task
              complete; reported as a percentage of tasks. \\
pp          & Percentage points. A move from 14\% to 17\% is $+3$ pp. \\
oracle & An idealised chooser told, after the fact, which modes solved a task,
allowed to pick one of them (cheapest where more than one worked). It cannot
be built; it measures how much opportunity exists. \\
ceiling & How much better that oracle does than the best single mode. An upper
bound on what any chooser could \emph{gain}, not a result any system attains. \\
drop-one contribution & How far the ceiling falls if one mode is deleted from
the oracle's menu. Running one mode a second time produces a number of the same
shape, so a positive value alone does not show irreplaceability
(Chapter~\ref{ch:repr}). \\
arm & One run of one mode over the whole task set. Six modes run once is six
arms; one mode run six times is also six arms — this dissertation matches the
count before comparing an arm's worth. \\
noise floor & How much a result moves when an unchanged setup is run again.
Two versions exist and must not be mixed: one for ceiling claims (what one more
run buys, 4.46--7.59\pp{}), one for mode-to-mode comparisons
(3.52--4.15\pp{}). Appendix~\ref{app:noise} derives both. \\
routing     & Choosing which mode to use per task, rather than fixing one in
              advance; \emph{always-cheapest} is the free fixed policy a router
              has to beat. \\
triage & The easier half of routing: predicting whether \emph{any} mode will
solve a task at all, sending it to the cheapest mode if not. \\
Pareto win  & Beating a comparison policy on \emph{both} axes at once: higher
              task success \emph{and} lower cost. \\
serving-time information & What a deployed system knows before it has paid for
              anything: task wording and cheap page text. Excludes anything
              the benchmark supplies that a deployment would not have, such as
              its own human-written difficulty rating. \\
nested cross-validation & A testing procedure in which every modelling choice
              is made on training tasks alone and judged on tasks held back,
              so a method cannot be tuned on the data it is scored on. \\
AUROC & How well a predictor \emph{ranks} one class above the other (0.5 =
chance, 1.0 = perfect). Says nothing about whether acting on that ranking
pays once the two kinds of mistake cost differently
(Chapters~\ref{ch:pred},~\ref{ch:e2e}). \\
\bottomrule
\end{tabularx}
\end{center}

    \thispagestyle{guide}
}
\newpage
